\chapter{Тензоры и релятивистская сила Лоренца}

Сегодня мы переходим от скалярных полей к электродинамике. Предыдущее обсуждение установило полезную схему: один функционал действия описывает и то, как поле действует на частицу, и то, как частица воздействует обратно на поле. В дальнейшем мы применим эту схему к электромагнитному полю, но сначала нам нужны обозначения, достаточно сильные, чтобы выявить симметрию вычисления.

Эта предварительная работа — не украшение. Знак минус, нуль, знак равенства и векторные обозначения — каждый из них сжимает идею, которую иначе пришлось бы повторять словами. Тензорные обозначения делают то же самое для теории относительности. Как только введены верхние и нижние индексы, свёртки и законы преобразования, электрические и магнитные силы сами соберутся в одно четырёхмерное уравнение.

\section{Четырёхвекторы и экономия индексов}

\subsection{Верхние и нижние компоненты}

Греческие индексы пробегают все четыре координаты пространства-времени, тогда как латинские — только три пространственные координаты:
\begin{equation}
\mu,\nu,\rho,\sigma=0,1,2,3,
\qquad
m,n=1,2,3.
\end{equation}
Произвольный контравариантный четырёхвектор имеет вид \(A^\mu=(A^0,A^1,A^2,A^3)\).
В используемом здесь соглашении
\begin{equation}
\eta_{\mu\nu}
=
\begin{pmatrix}
-1&0&0&0\\
0&1&0&0\\
0&0&1&0\\
0&0&0&1
\end{pmatrix},
\qquad
A_\mu=\eta_{\mu\nu}A^\nu .
\label{eq:sr6-metric-lowering}
\end{equation}
Следовательно, опускание индекса меняет знак временной компоненты и оставляет пространственные компоненты без изменения:
\begin{equation}
A_0=-A^0,
\qquad
A_m=A^m.
\label{eq:sr6-component-lowering}
\end{equation}
В плоском пространстве-времени этим эта операция и исчерпывается, хотя в общей теории относительности те же обозначения приобретут более богатый геометрический смысл.

\begin{figure}[htbp]
\centering
\includegraphics[width=0.90\textwidth]{lecture_06_figure_01.png}
\caption{Метрика Минковского и покомпонентное правило опускания индекса четырёхвектора.}
\lectureframe{00:07:45}
\end{figure}

Мы используем соглашение Эйнштейна о суммировании, когда индекс повторяется один раз сверху и один раз снизу. Поэтому свёртка
\begin{equation}
A_\mu B^\mu
=
-A^0B^0+A^1B^1+A^2B^2+A^3B^3
\label{eq:sr6-vector-contraction}
\end{equation}
является скаляром. Если сумма не имеет такой повторяющейся формы с верхним--нижним индексом, знак суммирования мы будем писать явно.

\subsection{Производная несёт нижний индекс}

Четыре оператора дифференцирования собираются в
\begin{equation}
\partial_\mu
\equiv
\frac{\partial}{\partial x^\mu}.
\end{equation}
Действуя на скалярное поле \(\phi(x)\), они образуют ковариантный вектор
\(\partial_\mu\phi\). Быстрее всего понять расположение индекса, если рассмотреть бесконечно малое смещение:
\begin{equation}
d\phi
=
\frac{\partial\phi}{\partial x^\mu}\,dx^\mu
=
\partial_\mu\phi\,dx^\mu.
\label{eq:sr6-total-differential}
\end{equation}
Левая часть — это разность значений скаляра в соседних событиях, следовательно, скаляр. Поскольку \(dx^\mu\) контравариантен, объект, свёрнутый с ним, должен быть ковариантным.

\begin{classroomqa}
\audiencequestion{Если у \(x^\mu\) верхний индекс, почему дифференцирование по \(x^\mu\) даёт нижний индекс?}
\lecturerresponse{Потому что дифференциальное изменение \(d\phi=\partial_\mu\phi\,dx^\mu\) является скаляром. Контравариантное смещение может дать этот скаляр только через свёртку с ковариантным градиентом. Поэтому положение индекса у производной определяется её законом преобразования, а не одной лишь типографикой дроби. \lecturetimestamp{00:10:10}}
\end{classroomqa}

\begin{figure}[htbp]
\centering
\includegraphics[width=0.90\textwidth]{lecture_06_figure_02.png}
\caption{Скалярный дифференциал \(d\phi\) показывает, что \(\partial_\mu\phi\) является ковариантным вектором.}
\lectureframe{00:11:05}
\end{figure}

То же правило даёт ещё один скаляр. Если \(B^\mu(x)\) — четырёхвекторное поле, то
\begin{equation}
\partial_\mu B^\mu
\label{eq:sr6-four-divergence}
\end{equation}
есть его четырёхдивергенция. Она объединяет производную по времени от временной компоненты с тремя пространственными производными от соответствующих пространственных компонент.

\begin{figure}[htbp]
\centering
\includegraphics[width=0.90\textwidth]{lecture_06_figure_03.png}
\caption{Две скалярные свёртки: \(A_\mu B^\mu\) и четырёхдивергенция \(\partial_\mu B^\mu(x)\).}
\lectureframe{00:14:00}
\end{figure}

У всего этого есть трёхмерный аналог. В евклидовом трёхмерном пространстве метрика является единичной матрицей, так что поднятие или опускание латинского индекса ничего не меняет. Поэтому мы будем поднимать и опускать пространственные индексы, когда это удобно, продолжая при этом тщательно различать индексы пространства-времени.

\section{Преобразования Лоренца в матричной форме}

Скаляры и векторы определяются тем, как они преобразуются. Скаляр сохраняет то же значение, тогда как контравариантный четырёхвектор подчиняется закону
\begin{equation}
(A^\mu)'=L^\mu{}_\nu A^\nu.
\label{eq:sr6-vector-transform}
\end{equation}
Видимый индекс \(\mu\) является свободным и присутствует по обеим сторонам. Повторяющийся индекс \(\nu\) суммируется и исчезает из окончательного покомпонентного уравнения.

Здесь под \emph{преобразованием Лоренца} понимаются как пространственные повороты, так и бусты. Буст может быть направлен вдоль любой оси; составляя повороты с бустами, можно получить общее преобразование. Пространственный поворот сам по себе не изменяет временную координату.

\begin{classroomqa}
\audiencequestion{Когда мы поворачиваем пространство, время остаётся полностью незатронутым?}
\lecturerresponse{Да. При обычном пространственном повороте временная компонента не изменяется. Такие повороты входят в группу Лоренца наряду с преобразованиями, которые смешивают время с пространством. \lecturetimestamp{00:16:27}}
\end{classroomqa}

Для буста со скоростью \(v\) вдоль оси \(x\), при \(c=1\),
\begin{equation}
L^\mu{}_\nu
=
\begin{pmatrix}
\gamma&-\gamma v&0&0\\
-\gamma v&\gamma&0&0\\
0&0&1&0\\
0&0&0&1
\end{pmatrix},
\qquad
\gamma=\frac{1}{\sqrt{1-v^2}}.
\label{eq:sr6-boost-matrix}
\end{equation}
Умножение на \((t,x,y,z)^{\mathsf T}\) даёт
\begin{equation}
\begin{aligned}
t'&=\gamma(t-vx),\\
x'&=\gamma(x-vt),\\
y'&=y,\\
z'&=z.
\end{aligned}
\end{equation}
По мере того как матрица дописывается на доске, аудитория подсказывает нули и два единичных элемента.

\begin{figure}[htbp]
\centering
\includegraphics[width=0.90\textwidth]{lecture_06_figure_04.png}
\caption{Матрица буста вдоль \(x\), действующая на столбец \((t,x,y,z)^{\mathsf T}\).}
\lectureframe{00:21:35}
\end{figure}

Поворот на угол \(\theta\) в плоскости \(y\)--\(z\), напротив, задаётся матрицей
\begin{equation}
R_{yz}(\theta)
=
\begin{pmatrix}
1&0&0&0\\
0&1&0&0\\
0&0&\cos\theta&\sin\theta\\
0&0&-\sin\theta&\cos\theta
\end{pmatrix}.
\label{eq:sr6-yz-rotation}
\end{equation}
Это означает
\begin{equation}
y'=y\cos\theta+z\sin\theta,
\qquad
z'=-y\sin\theta+z\cos\theta,
\end{equation}
тогда как \(t\) и \(x\) остаются неизменными. Произведения матриц, таких как
\eqref{eq:sr6-boost-matrix} и \eqref{eq:sr6-yz-rotation}, порождают более общие преобразования Лоренца.

\begin{figure}[htbp]
\centering
\includegraphics[width=0.90\textwidth]{lecture_06_figure_05.png}
\caption{Чистый поворот в плоскости \(y\)--\(z\), при котором \(t\) и \(x\) остаются без изменения.}
\lectureframe{00:23:10}
\end{figure}

Ковариантный вектор преобразуется матрицей с противоположным расположением индексов:
\begin{equation}
(A_\mu)'=M_\mu{}^\nu A_\nu.
\label{eq:sr6-covector-transform}
\end{equation}
Поправка из аудитории у доски замечает временно повторённый свободный индекс: индекс \(\mu\), по которому слева не ведётся суммирование, должен и справа быть индексом без суммирования. Эти две матрицы представляют одно и то же физическое изменение системы отсчёта и связаны соотношением
\begin{equation}
M=\eta L\eta,
\qquad
\eta^{-1}=\eta.
\label{eq:sr6-M-L-relation}
\end{equation}
Последнее тождество немедленно следует из того, что каждый диагональный элемент \(\eta\) совпадает со своим обратным.

\begin{figure}[htbp]
\centering
\includegraphics[width=0.90\textwidth]{lecture_06_figure_06.png}
\caption{Контравариантный и ковариантный законы преобразования, связанные соотношением \(M=\eta L\eta\).}
\lectureframe{00:25:20}
\end{figure}

Теперь центральная идея у нас в руках: четырёхвектор — это объект, компоненты которого преобразуются так же, как координаты пространства-времени.

\section{Тензоры: по одному преобразованию на каждый индекс}

\subsection{От произведений векторов к определению}

Скаляр — тензор нулевого ранга, вектор — тензор первого ранга, а объект с двумя индексами — тензор второго ранга. Начнём с двух векторов и образуем
\begin{equation}
T^{\mu\nu}=A^\mu B^\nu.
\label{eq:sr6-simple-tensor}
\end{equation}
Здесь \(4\times4=16\) компонент. Они идут от \(A^0B^0\) и
\(A^0B^1\) до \(A^3B^3\).
Символ \(T^{\mu\nu}\) обозначает весь этот массив, а не одно число.

Преобразуя каждый вектор по отдельности, получаем
\begin{align}
(A^\mu)'&=L^\mu{}_\sigma A^\sigma,\\
(B^\nu)'&=L^\nu{}_\tau B^\tau,
\end{align}
и, следовательно,
\begin{equation}
(T^{\mu\nu})'
=
L^\mu{}_\sigma L^\nu{}_\tau T^{\sigma\tau}.
\label{eq:sr6-rank-two-transform}
\end{equation}
Каждому тензорному индексу соответствует своя матрица Лоренца. Немые индексы
\(\sigma\) и \(\tau\) можно переименовывать как угодно, при условии, что каждая пара остаётся согласованной.

\begin{classroomqa}
\audiencequestion{Предполагает ли эта перестановка, что матрицы и компоненты векторов коммутируют?}
\lecturerresponse{Здесь речь идёт об обычных классических числах, а не о некоммутирующих квантовых операторах, поэтому они коммутируют. Более того, на оба индекса действует одна и та же матрица Лоренца. В этом вычислении вопроса о порядке не возникает. \lecturetimestamp{00:31:56}}
\end{classroomqa}

Это правило немедленно распространяется на третий ранг:
\begin{equation}
(T^{\mu\nu\lambda})'
=
L^\mu{}_\sigma
L^\nu{}_\tau
L^\lambda{}_\kappa
T^{\sigma\tau\kappa},
\label{eq:sr6-rank-three-transform}
\end{equation}
и аналогично — на любой ранг. Именно этот закон преобразования, а не некоторое конкретное разложение на векторы, и определяет тензор. Например,
\begin{equation}
A^\mu B^\nu+C^\mu D^\nu
\end{equation}
по-прежнему является тензором, хотя в общем случае его нельзя записать как одно произведение двух векторов.

\begin{figure}[htbp]
\centering
\includegraphics[width=0.90\textwidth]{lecture_06_figure_07.png}
\caption{Лестница на доске: от произведения векторов к законам преобразования тензоров второго и третьего рангов.}
\lectureframe{00:33:20}
\end{figure}

\begin{classroomqa}
\audiencequestion{Что означает покомпонентно умножение матриц на тензор второго ранга?}
\lecturerresponse{Зафиксируйте два выходных индекса и просуммируйте по двум скрытым индексам. Например,
\[
(T^{31})'=L^3{}_\sigma L^1{}_\tau T^{\sigma\tau}.
\]
Поскольку и \(\sigma\), и \(\tau\) принимают четыре значения, эта одна выходная компонента содержит 16 слагаемых. \lecturetimestamp{00:34:35}}
\end{classroomqa}

\begin{figure}[htbp]
\centering
\includegraphics[width=0.90\textwidth]{lecture_06_figure_08.png}
\caption{Развёртывание фиксированной компоненты \((T^{31})'\) делает явными два независимых суммирования.}
\lectureframe{00:35:55}
\end{figure}

Тензорные уравнения имеют решающее преимущество. Если два тензора имеют равные компоненты в одной системе отсчёта, закон преобразования делает их равными в любой системе отсчёта. Эквивалентно, если каждая компонента тензора обращается в нуль в одной системе отсчёта, то весь тензор обращается в нуль в любой системе отсчёта. Одной компоненты, обращающейся в нуль, недостаточно; должны обращаться в нуль все компоненты. Вот почему уравнение, записанное между тензорами, переносит свою истинность из одной инерциальной системы отсчёта в другую.

\subsection{Поднятие и опускание индексов тензора}

Индекс тензора опускается в точности так же, как индекс вектора:
\begin{equation}
T^\mu{}_\nu=T^{\mu\sigma}\eta_{\sigma\nu}.
\label{eq:sr6-tensor-lowering}
\end{equation}
Для диагональной метрики из \eqref{eq:sr6-metric-lowering} каждый опущенный временной индекс вносит один знак минус. Следовательно,
\begin{equation}
T_{00}=T^{00},
\qquad
T_{01}=-T^{01}.
\label{eq:sr6-tensor-component-signs}
\end{equation}
Первое равенство содержит два изменения знака; второе — только одно.

\begin{figure}[htbp]
\centering
\includegraphics[width=0.90\textwidth]{lecture_06_figure_09.png}
\caption{Примеры опускания двух тензорных индексов: два временных индекса дают два знака минус, тогда как один временной индекс даёт один.}
\lectureframe{00:41:00}
\end{figure}

\begin{classroomqa}
\audiencequestion{Почему полезно иметь и ковариантную, и контравариантную формы? Есть ли в этом что-то большее, чем компактная запись?}
\lecturerresponse{Для текущего вычисления непосредственная польза состоит в компактной записи свёрток и автоматическом учёте знаков Минковского. Между векторами и ковекторами есть более глубокое геометрическое различие, но его лучше развивать в общей теории относительности. В плоском пространстве-времени правила для индексов настолько эффективны, что обычно работают прямо с ними. \lecturetimestamp{00:41:26}}
\end{classroomqa}

\begin{classroomqa}
\audiencequestion{Всегда ли тензор имеет четыре компоненты, или это характерно только для теории относительности?}
\lecturerresponse{Диапазон каждого индекса задаётся размерностью пространства. Тензор первого ранга в четырёхмерном пространстве-времени имеет четыре компоненты; тензор второго ранга имеет шестнадцать. В трёх измерениях соответствующие числа равны трём и девяти. Само тензорное правило не зависит от размерности. \lecturetimestamp{00:42:55}}
\end{classroomqa}

\begin{classroomqa}
\audiencequestion{Можно ли описать ковариантные и контравариантные объекты в более геометрических терминах?}
\lecturerresponse{Да, у них есть геометрическая интерпретация, но если заняться ею сейчас, электродинамика отложится ещё дальше. Для этого курса достаточно правил преобразования и манипулирования индексами; геометрическое различие станет неизбежным в общей теории относительности. \lecturetimestamp{00:43:51}}
\end{classroomqa}

\section{Симметрия, антисимметрия и электромагнитный тензор}

Для произвольного тензора второго ранга перестановка индексов меняет компоненту:
\begin{equation}
T^{\mu\nu}\neq T^{\nu\mu}.
\end{equation}
Действительно, \(A^\mu B^\nu\) и \(A^\nu B^\mu\) в общем случае различны; конкретные компоненты \(A^0B^1\) и \(A^1B^0\) не обязаны совпадать.

Есть два важных частных случая. Симметричный тензор удовлетворяет
\begin{equation}
T^{\mu\nu}=T^{\nu\mu},
\end{equation}
а конкретным симметричным сочетанием является
\begin{equation}
A^\mu B^\nu+A^\nu B^\mu.
\end{equation}
Антисимметричный тензор удовлетворяет
\begin{equation}
F^{\mu\nu}=-F^{\nu\mu},
\label{eq:sr6-antisymmetry}
\end{equation}
как и сочетание
\begin{equation}
A^\mu B^\nu-A^\nu B^\mu.
\end{equation}
Если положить индексы равными в \eqref{eq:sr6-antisymmetry}, получаем
\begin{equation}
F^{00}=-F^{00},
\qquad
F^{00}=0,
\end{equation}
и тот же аргумент устраняет каждую диагональную компоненту. Следовательно, антисимметричный тензор \(4\times4\) имеет шесть независимых компонент.

\begin{figure}[htbp]
\centering
\includegraphics[width=0.90\textwidth]{lecture_06_figure_10.png}
\caption{На доске сопоставлены симметричные и антисимметричные тензоры второго ранга.}
\lectureframe{00:47:20}
\end{figure}

Этих шести компонент как раз достаточно для трёх компонент электрического поля и трёх компонент магнитного поля. При соглашении о знаках, записанном на доске,
\begin{equation}
F^{\mu\nu}
=
\begin{pmatrix}
0&-E_1&-E_2&-E_3\\
E_1&0&-B_3&B_2\\
E_2&B_3&0&-B_1\\
E_3&-B_2&B_1&0
\end{pmatrix}.
\label{eq:sr6-em-field-tensor}
\end{equation}
Отражение относительно диагонали меняет все знаки на противоположные. Электрическое и магнитное поля, следовательно, не являются независимыми релятивистскими объектами. Буст может смешивать их точно так же, как он смешивает пространство и время: поле, которое в одной системе отсчёта является чисто электрическим, в другой может иметь магнитную компоненту.

\begin{figure}[htbp]
\centering
\includegraphics[width=0.90\textwidth]{lecture_06_figure_11.png}
\caption{Полностью заполненная антисимметричная матрица, шесть независимых элементов которой задаются \(\mathbf E\) и \(\mathbf B\).}
\lectureframe{00:50:20}
\end{figure}

Доска, на которой записано \eqref{eq:sr6-em-field-tensor}, остаётся видимой, пока обсуждение переходит к динамике частицы. Теперь мы снова используем обозначение \(A_\mu\), но уже не для общего вектора из предыдущих примеров, а для электромагнитного четырёхпотенциала.

\section{Действие заряженной частицы}

Есть две возможные отправные точки. Можно изучать динамику самого поля, что приводит к уравнениям Максвелла, или изучать заряженную частицу, движущуюся в заданном поле, что приводит к силе Лоренца. Мы начнём с частицы.

При малой скорости нужно получить
\begin{equation}
m\mathbf a
=
e\left(\mathbf E+\mathbf v\times\mathbf B\right),
\label{eq:sr6-nonrel-lorentz-force}
\end{equation}
причём в варианте с гауссовыми единицами магнитный член умножается на множитель \(1/c\). Мы полагаем \(c=1\). Два слагаемых в
\eqref{eq:sr6-nonrel-lorentz-force} в трёхвекторной записи выглядят по-разному,
но теория относительности должна показать, что это части одного объекта.

\begin{figure}[htbp]
\centering
\includegraphics[width=0.90\textwidth]{lecture_06_figure_12.png}
\caption{Нерелятивистская сила Лоренца — результат, который должно воспроизвести инвариантное действие.}
\lectureframe{00:54:00}
\end{figure}

\subsection{Свободное движение и векторный потенциал}

Для свободной релятивистской частицы
\begin{equation}
S_{\mathrm{free}}
=
-m\int_1^2 d\tau
=
\int_{t_1}^{t_2}dt\,
\left[-m\sqrt{1-\dot{\mathbf x}^{\,2}}\right],
\label{eq:sr6-free-action}
\end{equation}
где
\begin{equation}
\dot{\mathbf x}^{\,2}
=
(\dot x^1)^2+(\dot x^2)^2+(\dot x^3)^2.
\end{equation}
Величина в квадратных скобках — это лагранжиан свободной частицы. Его импульс и энергия были найдены ранее; мы оставляем действие и спрашиваем, какой зависящий от поля член можно добавить.

\begin{figure}[htbp]
\centering
\includegraphics[width=0.90\textwidth]{lecture_06_figure_13.png}
\caption{Свободное действие как собственное время, накопленное вдоль мировой линии.}
\lectureframe{00:55:10}
\end{figure}

Основной электромагнитной переменной является четырёхвекторный потенциал \(A_\mu(x)\). Электрическое и магнитное поля будут выражены через него. На малом участке мировой линии \(dx^\mu\) свёртка
\begin{equation}
A_\mu(x)\,dx^\mu
\end{equation}
является скаляром. Суммирование этих скалярных приращений от одной конечной точки до другой даёт лоренц-инвариантное взаимодействие:
\begin{equation}
S_{\mathrm{int}}
=
-e\int_1^2 A_\mu(x)\,dx^\mu.
\label{eq:sr6-interaction-action}
\end{equation}
Коэффициент \(e\) — это электрический заряд. Общий знак связан с исторически сложившейся конвенцией для заряда; изменение конвенции меняет знак, приписываемый \(e\), а не содержание теории.

\begin{figure}[htbp]
\centering
\includegraphics[width=0.90\textwidth]{lecture_06_figure_14.png}
\caption{Участок мировой линии свёртывается с четырёхпотенциалом, образуя скалярное приращение \(A_\mu dx^\mu\).}
\lectureframe{00:59:20}
\end{figure}

\subsection{Форма в обычном времени}

Чтобы воспользоваться привычными уравнениями Эйлера--Лагранжа, разделим \eqref{eq:sr6-interaction-action} на \(dt\) и затем умножим на него. Поскольку
\begin{equation}
\frac{dx^0}{dt}=1,
\qquad
\frac{dx^m}{dt}=\dot x^m,
\end{equation}
полное действие имеет вид
\begin{equation}
S
=
\int dt\,
\left[
-m\sqrt{1-\dot{\mathbf x}^{\,2}}
-eA_0(x,t)
-e\dot x^m A_m(x,t)
\right].
\label{eq:sr6-complete-particle-action}
\end{equation}
Следовательно,
\begin{equation}
\boxed{
L
=
-m\sqrt{1-\dot{\mathbf x}^{\,2}}
-eA_0
-e\dot x^mA_m
}.
\label{eq:sr6-charged-lagrangian}
\end{equation}
Пока что \(A_\mu(x,t)\) — известное внешнее поле. Мы варьируем мировую линию,
а не поле.

\begin{figure}[htbp]
\centering
\includegraphics[width=0.90\textwidth]{lecture_06_figure_15.png}
\caption{Инвариантное взаимодействие вдоль мировой линии, переписанное в виде лагранжиана с обычным временем.}
\lectureframe{01:05:20}
\end{figure}

\section{Подробный вывод силы Лоренца}

Для пространственной компоненты \(x^m\) уравнение Эйлера--Лагранжа имеет вид
\begin{equation}
\frac{d}{dt}\frac{\partial L}{\partial\dot x^m}
=
\frac{\partial L}{\partial x^m}.
\label{eq:sr6-EL}
\end{equation}
Дифференцирование по скорости даёт
\begin{equation}
\frac{\partial L}{\partial\dot x^m}
=
\frac{m\dot x_m}{\sqrt{1-\dot{\mathbf x}^{\,2}}}
-eA_m.
\label{eq:sr6-canonical-momentum}
\end{equation}
Буква \(m\), к сожалению, обозначает и массу, и пространственный индекс. Контекст позволяет их различить: первое \(m\) в
\eqref{eq:sr6-canonical-momentum} — это масса; нижний индекс \(m\) обозначает одну из трёх пространственных компонент.

\begin{figure}[htbp]
\centering
\includegraphics[width=0.90\textwidth]{lecture_06_figure_16.png}
\caption{Производная импульса содержит и механический член, и член с векторным потенциалом.}
\lectureframe{01:09:20}
\end{figure}

\begin{classroomqa}
\audiencequestion{Не пропущен ли в первой записи левой части один член?}
\lecturerresponse{Да. Если дифференцировать только член свободной частицы, теряется зависимость от скорости у \(-e\dot x^mA_m\). Полная производная равна \(m\dot x_m/\sqrt{1-\dot{\mathbf x}^{\,2}}-eA_m\). Эта поправка существенна, потому что её производная по времени даёт часть и электрической, и магнитной силы. \lecturetimestamp{01:08:18}}
\end{classroomqa}

Правая часть \eqref{eq:sr6-EL} имеет вид
\begin{equation}
\frac{\partial L}{\partial x^m}
=
-e\,\partial_mA_0
-e\,\dot x^n\partial_mA_n.
\label{eq:sr6-position-derivative}
\end{equation}
Следовательно,
\begin{equation}
\frac{d}{dt}
\left(
\frac{m\dot x_m}{\sqrt{1-\dot{\mathbf x}^{\,2}}}
-eA_m
\right)
=
-e\,\partial_mA_0
-e\,\dot x^n\partial_mA_n.
\label{eq:sr6-EL-expanded}
\end{equation}

\begin{classroomqa}
\audiencequestion{Не были ли в этом вычислении перепутаны ковариантные и контравариантные пространственные компоненты?}
\lecturerresponse{Их действительно поднимали и опускали, но для латинских пространственных индексов евклидов пространственный блок метрики является единичным. Поэтому \(x_m=x^m\) и \(A_m=A^m\). Изменение знака возникает только у временного индекса. \lecturetimestamp{01:09:29}}
\end{classroomqa}

\begin{classroomqa}
\audiencequestion{Повторяющиеся латинские индексы тоже суммируются?}
\lecturerresponse{Да. По повторяющемуся латинскому индексу суммирование ведётся от \(1\) до \(3\); по повторяющемуся греческому индексу — от \(0\) до \(3\). Одиночный индекс, такой как \(m\), которым помечено одно пространственное уравнение, является свободным и не суммируется. Символ массы не является индексом. \lecturetimestamp{01:10:00}}
\end{classroomqa}

\begin{figure}[htbp]
\centering
\includegraphics[width=0.90\textwidth]{lecture_06_figure_17.png}
\caption{Уравнение Эйлера--Лагранжа до раскрытия полной производной по времени от векторного потенциала.}
\lectureframe{01:14:20}
\end{figure}

\subsection{Правило цепочки вдоль движущейся мировой линии}

Даже поле, не зависящее от времени, меняется вдоль траектории частицы, потому что частица проходит через разные точки. Полная производная равна
\begin{equation}
\frac{dA_m(x(t),t)}{dt}
=
\partial_0A_m+\dot x^n\partial_nA_m.
\label{eq:sr6-convective-derivative}
\end{equation}
Первый член — это явная зависимость от времени; второй — неявная зависимость от времени через \(x^n(t)\).

\begin{figure}[htbp]
\centering
\includegraphics[width=0.90\textwidth]{lecture_06_figure_18.png}
\caption{Векторный потенциал меняется явно со временем и неявно по мере того, как частица движется в пространстве.}
\lectureframe{01:17:20}
\end{figure}

Подставляя \eqref{eq:sr6-convective-derivative} в
\eqref{eq:sr6-EL-expanded} и перенося члены с полем вправо, получаем
\begin{equation}
\begin{aligned}
m\frac{d}{dt}
\left(
\frac{\dot x_m}{\sqrt{1-\dot{\mathbf x}^{\,2}}}
\right)
&=
e\left(\partial_0A_m-\partial_mA_0\right)
\\
&\quad +e\dot x^n
\left(\partial_nA_m-\partial_mA_n\right).
\end{aligned}
\label{eq:sr6-grouped-force}
\end{equation}
В вычислении наконец выявился организующий принцип. Первые скобки не содержат скорости; вторые линейны по скорости. Это в точности соответствует разделению на электрическую и магнитную силы.

\begin{figure}[htbp]
\centering
\includegraphics[width=0.90\textwidth]{lecture_06_figure_19.png}
\caption{Силовые члены сгруппированы в не зависящую от скорости электрическую часть и зависящую от скорости антисимметричную часть.}
\lectureframe{01:23:40}
\end{figure}

Определим электрические компоненты в соглашении, принятом в этом выводе, как
\begin{equation}
E_m=\partial_0A_m-\partial_mA_0.
\label{eq:sr6-electric-from-potential}
\end{equation}
Пространственная антисимметричная комбинация
\begin{equation}
C_{mn}
\equiv
\partial_nA_m-\partial_mA_n
\label{eq:sr6-spatial-curl-tensor}
\end{equation}
содержит три независимые компоненты ротора векторного потенциала. Ее ориентация выбрана так, что
\begin{equation}
\dot x^nC_{mn}
=
\left(\mathbf v\times\mathbf B\right)_m.
\end{equation}
Следовательно, уравнение \eqref{eq:sr6-grouped-force} представляет собой релятивистскую форму по импульсу для \eqref{eq:sr6-nonrel-lorentz-force}. Величина под производной по времени — это релятивистский импульс на единицу массы, а не одна лишь ньютоновская скорость.

Этот вывод стоит повторить вручную: начните с \eqref{eq:sr6-charged-lagrangian}, вычислите обе части \eqref{eq:sr6-EL}, примените правило цепочки и сгруппируйте члены по их зависимости от скорости. Это упражнение одновременно служит практикой по вариационной механике, многомерному анализу и электродинамике.

\begin{classroomqa}
\audiencequestion{Почему в качестве электромагнитного взаимодействия выбирают \(-e\int A_\mu dx^\mu\)?}
\lecturerresponse{В конечном счете форму связи определяет эксперимент, но два структурных требования делают это естественным первым кандидатом. Действие составляется из локальных приращений вдоль траектории, и каждое приращение должно быть скаляром. Скаляр \(dx^\mu dx_\mu=d\tau^2\) уже представлен свободным действием. Свёртка \(dx^\mu\) с полем \(A_\mu\) дает самый простой новый зависящий от поля скаляр. Можно записать и более сложные скаляры, а калибровочная инвариантность позднее сделает выбор более определенным. \lecturetimestamp{01:25:34}}
\end{classroomqa}

\begin{classroomqa}
\audiencequestion{Можно ли также вывести уравнения Максвелла из принципа наименьшего действия?}
\lecturerresponse{Да. Настоящее действие описывает частицу в заданном электромагнитном поле. Чтобы получить уравнения Максвелла, нужно добавить собственное действие поля и варьировать само поле, как мы делали для скалярного поля. \lecturetimestamp{01:28:58}}
\end{classroomqa}

\begin{classroomqa}
\audiencequestion{Где в сгруппированном уравнении для силы скрыто суммирование?}
\lecturerresponse{Свободный индекс \(m\) обозначает одно из трех компонентных уравнений и не суммируется. Повторяющийся \(n\) в \(\dot x^n(\partial_nA_m-\partial_mA_n)\) суммируется по \(1,2,3\). Если записать этот член без обозначений Эйнштейна, в нем явно стояло бы \(\sum_{n=1}^{3}\). \lecturetimestamp{01:29:26}}
\end{classroomqa}

\begin{classroomqa}
\audiencequestion{Почему отождествлять \(\partial_0A_m-\partial_mA_0\) с электрическим полем, а не предполагать заранее обычную формулу?}
\lecturerresponse{Действие дало член силы, который зависит от положения и времени, но не от скорости частицы. Именно такова операциональная роль электрического поля. Оставшиеся антисимметричные члены умножаются на скорость и поэтому играют магнитную роль. На это вычисление можно смотреть как на обнаружение полей из закона силы, а не как на подстановку отдельных определений заранее. \lecturetimestamp{01:31:05}}
\end{classroomqa}

\section{Явно ковариантный закон силы}

Левая часть \eqref{eq:sr6-grouped-force} пока еще не записана в явно четырехмерной форме. Введем четырехскорость
\begin{equation}
u^\mu=\frac{dx^\mu}{d\tau}.
\label{eq:sr6-four-velocity}
\end{equation}
Ее пространственные компоненты удовлетворяют
\begin{equation}
u^m
=
\frac{dx^m}{d\tau}
=
\frac{\dot x^m}{\sqrt{1-\dot{\mathbf x}^{\,2}}}.
\end{equation}
Умножение пространственного уравнения на \(dt/d\tau\) переводит
\(d/dt\) в \(d/d\tau\):
\begin{equation}
\frac{dt}{d\tau}\frac{d}{dt}
=
\frac{d}{d\tau}.
\end{equation}
Следовательно, левая часть становится пространственной частью собственного ускорения,
\begin{equation}
m\frac{d^2x^m}{d\tau^2}.
\end{equation}

Правую часть можно записать, не разделяя время и пространство:
\begin{equation}
m\frac{d^2x^m}{d\tau^2}
=
e\left(\partial_\nu A_m-\partial_mA_\nu\right)
\frac{dx^\nu}{d\tau}.
\label{eq:sr6-spatial-covariant-force}
\end{equation}
При \(\nu=0\) это воспроизводит электрическую скобку; при \(\nu=1,2,3\) — магнитную.

Теперь обе части имеют закон преобразования четырехвектора. Следовательно, три пространственных уравнения являются компонентами
\begin{equation}
\boxed{
m\frac{d^2x^\mu}{d\tau^2}
=
e\left(\partial_\nu A^\mu-\partial^\mu A_\nu\right)u^\nu
}.
\label{eq:sr6-covariant-force}
\end{equation}
Если определить
\begin{equation}
F_{\mu\nu}
\equiv
\partial_\mu A_\nu-\partial_\nu A_\mu,
\label{eq:sr6-F-from-A}
\end{equation}
то порядок индексов и знак взаимодействия, использованные в этой лекции, дают эквивалентную компактную форму
\begin{equation}
m\frac{du^\mu}{d\tau}
=
-eF^\mu{}_\nu u^\nu.
\label{eq:sr6-covariant-force-F}
\end{equation}
\editorialnote{В учебниках по электродинамике знаки по-разному распределяются между сигнатурой метрики, определением \(A_\mu\), зарядом и \(F_{\mu\nu}\). Уравнения \eqref{eq:sr6-covariant-force} и \eqref{eq:sr6-covariant-force-F} сохраняют знак и порядок индексов, полученные из действия \(-e\int A_\mu dx^\mu\); инвариантным содержанием является антисимметризованная производная потенциала.}

\begin{figure}[htbp]
\centering
\includegraphics[width=0.90\textwidth]{lecture_06_figure_20.png}
\caption{Три пространственных уравнения, собранные в четырехвекторное уравнение с использованием собственного времени.}
\lectureframe{01:42:40}
\end{figure}

\begin{classroomqa}
\audiencequestion{Почему допустимо добавить четвертое уравнение — для временной компоненты? Разве \(d^2x^0/d\tau^2\) просто не равно нулю?}
\lecturerresponse{Вообще говоря, оно не равно нулю: \(x^0=t\), но \(dt/d\tau=\gamma\) меняется, когда меняется скорость частицы. Решающий аргумент — ковариантность. Скалярное действие гарантирует лоренц-инвариантные уравнения; как только пространственные компоненты распознаются как компоненты одного четырехвекторного уравнения, преобразования Лоренца смешивают их с временной компонентой. Четвертое уравнение должно сопровождать остальные три. \lecturetimestamp{01:42:40}}
\end{classroomqa}

\subsection{Что означает четвертая компонента}

Три пространственные компоненты выражают релятивистское утверждение
\begin{equation}
\frac{d\mathbf p}{dt}=\mathbf F.
\end{equation}
Энергия и импульс сами являются компонентами четырехвектора, поэтому временная компонента должна описывать соответствующий баланс энергии. На языке трехмерных векторов,
\begin{equation}
\frac{dK}{dt}=\mathbf F\cdot\mathbf v.
\label{eq:sr6-work-energy}
\end{equation}
Оно задает скорость, с которой поле совершает работу над частицей. Магнитный вклад перпендикулярен \(\mathbf v\) и работы не совершает; электрический вклад изменяет энергию частицы.

\begin{classroomqa}
\audiencequestion{Если ковариантность формально дает четвертое уравнение, как мы узнаем, что оно имеет физический смысл?}
\lecturerresponse{Первые три компоненты говорят, что сила изменяет импульс. Четвертая говорит, что работа изменяет энергию: \(dK/dt=\mathbf F\cdot\mathbf v\). Поскольку энергия и импульс образуют один четырехвектор, их отклик на поле тоже должен представлять собой одно четырехкомпонентное утверждение. \lecturetimestamp{01:45:20}}
\end{classroomqa}

\begin{figure}[htbp]
\centering
\includegraphics[width=0.90\textwidth]{lecture_06_figure_21.png}
\caption{Пространственный закон \(d\mathbf p/dt=\mathbf F\) и соответствующее ему уравнение для временной компоненты, соотношение работа--энергия.}
\lectureframe{01:47:10}
\end{figure}

Метод важнее окончательного сжатия записи. Мы потребовали, чтобы действие было скаляром еще до варьирования. Получающиеся уравнения затем можно было перегруппировать в тензоры, не исправляя их симметрию задним числом. В современной теории поля это превращается в общее правило построения: конструируйте лагранжиан так, чтобы он согласовывался с симметриями задачи.

\section{Локальность, инвариантность Лоренца и следующая симметрия}

Три принципа вновь и вновь возникают во всей релятивистской теории поля. Два уже были использованы; третий будет развит дальше.

\subsection{Локальность}

Принцип локальности утверждает, что изменение частицы или поля определяется тем, что происходит бесконечно близко в пространстве-времени. Частица откликается на поле, вычисленное на ее мировой линии, а не на несвязанное с ней значение где-то далеко. Поле в одной точке откликается на поле и его производные в окрестности этой точки, а не непосредственно на произвольные удаленные значения или на будущие данные.

Для поля действие поэтому принимает локальную форму
\begin{equation}
S
=
\int d^3x\,dt\,
\mathcal L\!\left(\phi,\partial_\mu\phi\right).
\label{eq:sr6-local-action}
\end{equation}
Интеграл складывает по одному вкладу от каждой точки пространства-времени. Первые производные кодируют сравнение с соседними точками, а вариационный принцип приводит к дифференциальным уравнениям. Можно было бы представить себе нелокальные произведения, связывающие поля в широко разделенных точках, но тогда их уравнения делали бы отклик здесь напрямую зависящим от условий там. Это не та локальная структура, которая здесь вводится. Действия с высшими производными возможны и образуют интересный предмет, но для настоящего построения они не нужны.

\begin{figure}[htbp]
\centering
\includegraphics[width=0.90\textwidth]{lecture_06_figure_22.png}
\caption{Локальность, выраженная как пространственно-временной интеграл от плотности, построенной из полей и их локальных производных.}
\lectureframe{01:51:00}
\end{figure}

\begin{classroomqa}
\audiencequestion{Представляет ли \(\partial_\mu\phi\) четыре компоненты производной?}
\lecturerresponse{Да. Это обозначает производную по времени и три пространственные производные \(\phi\), собранные в один ковариантный четырехвектор. \lecturetimestamp{01:52:25}}
\end{classroomqa}

\begin{classroomqa}
\audiencequestion{Как в принципе локальности количественно определяется \emph{близость}?}
\lecturerresponse{С помощью производных. Первая производная сравнивает поле в точке с его значением в бесконечно близких соседних точках. Локальный лагранжиан зависит от полей и конечного набора таких производных в одной и той же точке пространства-времени. \lecturetimestamp{01:52:55}}
\end{classroomqa}

\begin{classroomqa}
\audiencequestion{Как полевая версия локальности переводится на расчет для частицы, который мы только что завершили?}
\lecturerresponse{Лагранжиан частицы локален вдоль мировой линии: он зависит от соседних положений на мировой линии через скорость и от электромагнитного потенциала, взятого в текущем положении частицы. Ускорение частицы здесь не зависит непосредственно от поля в удаленной точке. \lecturetimestamp{01:53:31}}
\end{classroomqa}

\subsection{Инвариантность Лоренца}

Инвариантность Лоренца реализуется построением лагранжиана из скаляров. Наблюдатели могут разлагать эти скаляры на различные временные и пространственные компоненты, но относительно действия они согласны. Поэтому тензорные уравнения, полученные варьированием, имеют один и тот же вид в любой инерциальной системе отсчёта.

\subsection{Калибровочная инвариантность}

Калибровочная инвариантность — последний оставшийся принцип. Здесь он заявлен, но ещё не определён. На следующем шаге он наложит ограничение на то, какие свойства четырёхпотенциала физически значимы, и будет направлять выбор действия поля, которое приводит к уравнениям Максвелла.

\begin{figure}[htbp]
\centering
\includegraphics[width=0.90\textwidth]{lecture_06_figure_23.png}
\caption{Три организующих принципа: локальность, инвариантность Лоренца и калибровочная инвариантность.}
\lectureframe{01:55:40}
\end{figure}

Эти принципы вновь возникают в общей теории относительности, квантовой электродинамике, теории Янга--Миллса и Стандартной модели. Детали, такие как число видов частиц, меняются от теории к теории; более глубокую основу задают локальность и симметрия. Практическое задание поэтому состоит в том, чтобы заново проследить вычисление Эйлера--Лагранжа, пока не удастся объяснить каждый индекс и каждый сгруппированный член. Затем калибровочная инвариантность войдёт как следующее ограничение на то же самое электромагнитное действие.